\subsection{Adam}

\begin{frame}{Adam: Explicación Gráfica}
    \begin{block}{Concepto principal}
        [Explica la combinación de estimación de primer momento (inercia) y segundo momento (escala adaptativa)].
    \end{block}
    \begin{itemize}
        \item {[Idea gráfica 1: Corrección de sesgo durante las primeras iteraciones cerca del origen].}
        \item {[Idea gráfica 2: Trayectoria en superficies con diferente curvatura].}
    \end{itemize}
\end{frame}

\begin{frame}{Adam: Pseudocódigo}
    \begin{algorithmic}[1]
        \State \textbf{Input:} {[Parámetros \(\alpha, \beta_1, \beta_2, \epsilon\), punto inicial \(x\)]}
        \State Inicializar momentos \(m \leftarrow 0\), \(v \leftarrow 0\), paso \(t \leftarrow 0\)
        \While{[Condición de parada]}
            \State \(t \leftarrow t + 1\)
            \State Actualizar momentos sesgados \(m\) y \(v\) con el gradiente \(g\)
            \State Aplicar corrección de sesgo: \(\hat{m} \leftarrow m / (1-\beta_1^t)\), \(\hat{v} \leftarrow v / (1-\beta_2^t)\)
            \State Actualizar parámetros: \(x \leftarrow x - \alpha \hat{m} / (\sqrt{\hat{v}} + \epsilon)\)
        \EndWhile
        \State \textbf{Return} {[Punto óptimo]}
    \end{algorithmic}
\end{frame}

\begin{frame}{Adam: Ejemplo Paso a Paso}
    \begin{exampleblock}{Planteamiento del problema}
        [Define función cuadrática o de prueba junto con hiperparámetros estándar].
    \end{exampleblock}
    \begin{itemize}
        \item \textbf{Paso 1:} {[Cálculo de \(m_1, v_1\) y aplicación de corrección de sesgo].}
        \item \textbf{Paso 2:} {[Cálculo de la primera actualización de \(x\)].}
        \item \textbf{Paso 3:} {[Cálculos para la iteración \(t=2\)].}
    \end{itemize}
\end{frame}

\begin{frame}[fragile]{Adam: Código}
\begin{lstlisting}[language=Python]
def adam(grad_f, x0, alpha=0.001, beta1=0.9, beta2=0.999, eps=1e-8, max_iter=100):
    # TODO: Inicializar m, v y contador t
    # for t in range(1, max_iter + 1):
    #     g = grad_f(x)
    #     m = beta1 * m + (1 - beta1) * g
    #     v = beta2 * v + (1 - beta2) * (g ** 2)
    #     m_hat = m / (1 - beta1**t)
    #     v_hat = v / (1 - beta2**t)
    #     x = x - alpha * m_hat / (np.sqrt(v_hat) + eps)
    pass
\end{lstlisting}
\end{frame}

\begin{frame}{Adam: Análisis de Convergencia}
    \begin{alertblock}{Propiedades de Convergencia}
        [Resumir comportamiento general en problemas estocásticos y convexos].
    \end{alertblock}
    \begin{itemize}
        \item {[Impacto de la corrección de sesgo en los primeros pasos].}
        \item {[Discusión sobre casos donde puede no converger (variantes AMSGrad)].}
    \end{itemize}
\end{frame}